\documentclass[10pt,a4paper]{amsart}
\usepackage[margin=19mm]{geometry}
\usepackage{amsmath,amssymb,mathtools}
\usepackage[scr=rsfs,cal=euler]{mathalfa}
\usepackage{xfrac}
\usepackage[unicode,hidelinks,bookmarks=false]{hyperref}
\newcommand{\ie}{i.e.\ }
\newcommand{\cf}{cf.\ }
\newcommand{\resp}{respectively\ }
\newcommand{\Subsection}{\subsection}
\providecommand{\nobreakdash}{\nobreak\hbox{-}}
\setlength{\emergencystretch}{2em}
\multlinegap0pt
\usepackage{amssymb}
\usepackage{tikz}
\newtheorem{lemma}{Lemma}
\newcommand{\RV}{\textup{RV}}
\newcommand{\op}{\operatorname}
\title{Nonclassical Weyl Laws and\\ Connes' Integration for weak Lorentz Ideals, II}
\author{Raphael Ponge, Yongqiang Tian}
\date{Source: arXiv:2609.09568v1 (CC BY 4.0)}
\begin{document}
\maketitle

\noindent\emph{Source and licence.} Nonclassical Weyl Laws and\\ Connes' Integration for weak Lorentz Ideals, II. Version arXiv:2609.09568v1 (CC BY 4.0), distributed by arXiv under the Creative Commons Attribution 4.0 International licence, http://creativecommons.org/licenses/by/4.0/, as recorded in the arXiv licence metadata for this exact version and shown on the abstract page https://arxiv.org/abs/2609.09568v1. Full licence notice: \texttt{licenses/arxiv-2609.09568-CC-BY-4.0.txt}.\par\smallskip

\noindent\textit{Editorial scope.} Counted unit: the article's lemma lem:app.Ntolambda1 (source0.tex lines 1884--1893). The article's complete original proof of it is delivered with it (source0.tex lines 1894--1982, one \texttt{proof} environment of 89 lines). No mathematics is written, repaired or re-derived: the statement and the proof are the article's own lines, in the article's own notation, and no retained line is substituted or re-flowed.  Retained uncounted as the source-local context the proof needs: lines 1871--1873; lines 1876--1878.  Nothing was omitted from any retained span, so the package carries no omission record.  No retained line contains a citation, so the wrapper carries no bibliography and no bibliography entry.  No retained line contains a figure environment, an image-inclusion command or drawing code, so no image is delivered and the package declares no asset dependency.  Editorial formatting only, in addition: an AMS \texttt{amsart} preamble that restates the article's own declarations and macros used by the retained lines, namely the environments \texttt{lemma} on separate counters, together with the standard packages those lines need.  The retained environments print in this document as Lemma 1 in order of appearance, whereas the article prints the same items with the numbering of its own full document.  The complete native source of the article as distributed in the arXiv e-print for this exact version is bundled unchanged as \texttt{sources/209830/source0.tex}.
 \begin{equation}\label{cond:ii}
 h(at)=a^p h(t)\left(1+\op{O}(\epsilon(t))\right) \qquad \textup{as}\ t\rightarrow \infty.
\end{equation}

\begin{equation}\label{cond:hsharp}
 h^\sharp\circ h(t)= t\left(1+\op{O}(\epsilon(t))\right).
\end{equation}

\begin{lemma}\label{lem:app.Ntolambda1}
Assume that the conditions (i)--(iii) are satisfied, and as $ \lambda\rightarrow 0^+$ we have
\begin{equation*}
 N(\lambda)= h(\lambda^{-1}) \left[ c+\op{O}\left(\epsilon(\lambda^{-1})\right)\right]  \qquad \textup{with}\ c>0. 
\end{equation*}
Then,  we have
\begin{equation}\label{eq:case c>0 result}
 \lambda_j=c^{\frac1p} h^\sharp(j)^{-1} \left[ 1+\op{O}\left((\epsilon \circ h^\sharp)(j)\right)\right] \qquad \textup{as}\ j\rightarrow \infty.
\end{equation}
\end{lemma}

\begin{proof}
Let us first assume that $c=1$. Let $h^\natural:(0,\infty)\to(0,\infty)$ be a continuous function inverting $h(t)$ on some interval $[t_0,\infty)$. This is an $\RV_{1/p}$-function.
Set \begin{equation}\label{eq:case c=1 Nlambda}
N(\lambda)=h(\lambda^{-1})[1+\hat{\varepsilon}(\lambda^{-1})].
\end{equation} Note that  $|\hat{\varepsilon}(t)|=\op{O}(\epsilon(t))$. In addition, let $(a_j)_{j\geq1}\subset(0,\infty)$ be such that
\begin{equation}\label{eq:case c=1 a_j}
a_j=1-\epsilon(\lambda_j^{-1})
\end{equation} for $j$ large enough. Combining ~(\ref{eq:case c=1 Nlambda})--(\ref{eq:case c=1 a_j}) we get
\begin{equation*}
h(\lambda_j^{-1})(1-|\hat{\varepsilon}(\lambda_j^{-1})|)\leq N(\lambda_j)\leq j \leq N(a_j\lambda_j)\leq h((a_j\lambda_j)^{-1})(1+|\hat{\varepsilon}((a_j\lambda_j)^{-1})|).
\end{equation*}
As $h^\natural$ is ultimately increasing, for $j$ large enough, we have
\begin{equation*}
h^\natural[	h(\lambda_j^{-1})(1-|\hat{\varepsilon}(\lambda_j^{-1})|)]\leq h^\natural(j)\leq h^\natural[  h((a_j\lambda_j)^{-1})(1+|\hat{\varepsilon}((a_j\lambda_j)^{-1})|)].
\end{equation*}
Recall that, by assumption there is $q>p$ such that $t^{-q}h(t)$ is ultimately decreasing, i.e., there is $t_0>0$ such that
\begin{equation*}
\lambda^{-q}h(\lambda)\leq t^{-q}h(t),\qquad\mbox{for all}\quad t_0\leq \lambda\leq t.
\end{equation*}
For $\lambda=h^\natural[	h(\lambda_j^{-1})(1-|\hat{\varepsilon}(\lambda_j^{-1})|)]$ and $t=\lambda_j^{-1}$, we have
\begin{equation*}
h^\natural[	h(\lambda_j^{-1})(1-|\hat{\varepsilon}(\lambda_j^{-1})|)]^{-q}h(\lambda_j^{-1})\leq \lambda_j^qh(\lambda_j^{-1}).
\end{equation*}
That is,
\begin{equation}\label{eq:case c=1 lower}
\lambda_j^{-1}(1-|\hat{\varepsilon}(\lambda_j^{-1})|)^{1/q}\leq h^\natural[	h(\lambda_j^{-1})(1-|\hat{\varepsilon}(\lambda_j^{-1})|)].
\end{equation}
For $\lambda=(a_j\lambda_j)^{-1}$ and $t=h^\natural[  h((a_j\lambda_j)^{-1})(1+|\hat{\varepsilon}((a_j\lambda_j)^{-1})|)]$, we get
\begin{equation*}
(a_j\lambda_j)^qh((a_j\lambda_j)^{-1})\leq h^\natural[  h((a_j\lambda_j)^{-1})(1+|\hat{\varepsilon}((a_j\lambda_j)^{-1})|)]^{-q}h((a_j\lambda_j)^{-1})(1+|\hat{\varepsilon}((a_j\lambda_j)^{-1})|).
\end{equation*}
That is,
\begin{equation*}
h^\natural[  h((a_j\lambda_j)^{-1})(1+|\hat{\varepsilon}((a_j\lambda_j)^{-1})|)]\leq (a_j\lambda_j)^{-1}(1+|\hat{\varepsilon}((a_j\lambda_j)^{-1})|)^{1/q}.
\end{equation*}
Combining this with ~(\ref{eq:case c=1 lower}), we get
\begin{equation}\label{eq:case c=1 a}
\lambda_j^{-1}(1-|\hat{\varepsilon}(\lambda_j^{-1})|)^{1/q}\leq h^\natural(j)\leq a_j^{-1}\lambda_j^{-1}(1+|\hat{\varepsilon}((a_j\lambda_j)^{-1})|)^{1/q}.
\end{equation}
Clearly, we have $1-|\hat{\varepsilon}(\lambda_j^{-1})|=1-\op{O}(|\hat{\varepsilon}(\lambda_j^{-1})|)=1-\op{O}(\epsilon(\lambda_j^{-1}))$. As  $\epsilon$ is an $\RV_{\rho}$-function and $a_j\to1$, we have $ \epsilon((a_j\lambda_j)^{-1})\sim a_j^{-\rho}\epsilon(\lambda_j^{-1})\sim \epsilon(\lambda_j^{-1})$. Similarly, for $j$ large enough, we see from ~(\ref{eq:case c=1 a_j}) that
\begin{equation*}
a_j^{-1}(1+|\hat{\varepsilon}((a_j\lambda_j)^{-1})|)^{1/q}=(1-\epsilon(\lambda_j^{-1}))^{-1}[1+\op{O}(\epsilon((a_j\lambda_j)^{-1}))]^{1/q}=1+\op{O}(\epsilon(\lambda_j^{-1})).
\end{equation*}
Combining this with ~(\ref{eq:case c=1 a}) shows that
\begin{equation}\label{eq:case c=1 b}
\lambda_j=h^\natural(j)^{-1}(1+\op{O}(\epsilon(\lambda_j^{-1}))).
\end{equation}
It remains to prove~(\ref{eq:case c>0 result}) for $h^\sharp$, instead of $h^\natural$. In view of ~(\ref{cond:hsharp}), we have
\begin{equation*}
h^\sharp(t)=h^\sharp\circ h(h^\natural(t))=h^\natural(t)(1+\op{O}(\epsilon\circ h^\natural(t))).
\end{equation*}
Thus,
\begin{equation}\label{eq:case c=1 c}
h^{\natural}(t)^{-1}=h^\sharp(t)^{-1}(1+\op{O}(\epsilon\circ h^\natural(t))).
\end{equation}
This implies
$\lambda_j^{-1}\sim h^\sharp(j)$.
As $\epsilon$ is $\RV_\rho$, it follows that $\epsilon(\lambda_j^{-1})\sim \epsilon\circ h^\natural(j)$, and hence we have
\begin{equation*}
\lambda_j=h^\natural(j)^{-1}(1+\op{O}(\epsilon\circ h^\natural(j))).
\end{equation*}
In particular, $h^\sharp(t)\sim h^\natural(t)$. As $\epsilon$ is $\RV_\rho$, it follows that $\epsilon\circ h^\sharp(t)\sim \epsilon\circ h^\natural(t)$. Combining this with ~(\ref{eq:case c=1 b})--(\ref{eq:case c=1 c}), we have
\begin{equation*}
\lambda_j=h^\sharp(j)^{-1}(1+\op{O}(\epsilon\circ h^\sharp(j)))(1+\op{O}(\epsilon\circ h^\sharp(j)))=h^\sharp(j)^{-1}(1+\op{O}(\epsilon\circ h^\sharp(j))).
\end{equation*}
This proves the result $c=1$. 

Suppose now that $c$ is any positive constant. If we apply~(\ref{cond:ii}) to $a=c^{1/p}$ we get
\begin{equation*}
h(c^{1/p}t)=c h(t)\left(1+\op{O}(\epsilon(t))\right) \qquad \textup{as}\ t\rightarrow \infty.
\end{equation*}
Thus, if we set $h_c(t):=h(c^{1/p}t)$, then 
\begin{equation*}
N(\lambda)=h(\lambda^{-1}) \left[ c+\op{O}\left(\epsilon(\lambda^{-1})\right)\right]=ch(\lambda^{-1}) \left[ 1+\op{O}\left(\epsilon(\lambda^{-1})\right)\right]=h_c(\lambda^{-1}) \left[ 1+\op{O}\left(\epsilon(\lambda^{-1})\right)\right].
\end{equation*}
Setting $h_c^\sharp(t):=c^{-1/p}h^\sharp(t)$, we have
\begin{equation*}
h_c^\sharp\circ h_c(t)=c^{-1/p}h^\sharp \circ h(c^{1/p}t)=c^{-1/p}\cdot c^{1/p}t\cdot[1+\op{O}(\epsilon(c^{1/p}t))].
\end{equation*}
As $\epsilon$ is an $\RV_\rho$-function, we have $ \epsilon(c^{1/p}t)\sim c^{\rho/p}\epsilon(t)$, and so we obtain
\begin{equation*}
h_c^\sharp\circ h_c(t)=t(1+\op{O}(\epsilon(t))).
\end{equation*}
The result for $c=1$ then gives
\begin{equation*}
\lambda_j= h_c^\sharp(j)^{-1}(1+\op{O}(\epsilon(\lambda_j^{-1})))= c^{1/p}h^\sharp(j)^{-1}(1+\op{O}(\epsilon(\lambda_j^{-1})))=h^\sharp(j)^{-1}(c^{1/p}+\op{O}(\epsilon(\lambda_j^{-1}))).
\end{equation*}
This implies that $\lambda_j^{-1}\sim c^{-1/p}h^\sharp(j)$. As $\epsilon$ is $\RV_\rho$, it follows that $\epsilon(\lambda_j^{-1})\sim c^{-\rho/p}\epsilon \circ h^\sharp(j)$ and so we get the asymptotic~(\ref{eq:case c>0 result}). The proof is complete. 
\end{proof}

\end{document}
